\documentclass[10pt,aspectratio=169]{beamer}
\input{../../_en_preambulo_beamer}% Comandos y macros estadísticas compartidas para las presentaciones Beamer en inglés (Metropolis)

% Conjuntos numéricos básicos
\newcommand{\R}{\mathbb{R}}
\newcommand{\N}{\mathbb{N}}
\newcommand{\Q}{\mathbb{Q}}
\newcommand{\Z}{\mathbb{Z}}
\newcommand{\set}[1]{\left\{ #1 \right\}}
\newcommand{\abs}[1]{\left|#1\right|}
\newcommand{\norm}[1]{\left\| #1 \right\|}

% Comandos estadísticos y probabilísticos del curso
\newcommand{\Var}{\operatorname{Var}}
\newcommand{\cov}{\operatorname{Cov}}
\newcommand{\corr}{\rho}
\newcommand{\comb}[2]{\begin{pmatrix} #1 \\ #2 \end{pmatrix}}
\newcommand{\s}{\sigma}
\newcommand{\particion}{\mathfrak{P}}
\newcommand{\card}[1]{\#\left( #1 \right)}
\newcommand{\seq}[3]{\set{#1}_{#2}^{#3}}
\newcommand{\E}{\mathbb{E}}
\newcommand{\Prob}{\mathbb{P}}

% Entornos pedagógicos y cajas en inglés académico natural (soportando tanto nombres en inglés como en español)
\newenvironment{discussion}{%
  \begin{alertblock}{Discussion Question}%
}{%
  \end{alertblock}%
}
\newenvironment{discusion}{%
  \begin{alertblock}{Discussion Question}%
}{%
  \end{alertblock}%
}

\newenvironment{reflection}{%
  \begin{alertblock}{Classroom Reflection}%
}{%
  \end{alertblock}%
}
\newenvironment{reflexion}{%
  \begin{alertblock}{Classroom Reflection}%
}{%
  \end{alertblock}%
}

\newenvironment{paradox}{%
  \begin{alertblock}{Exploratory Paradox}%
}{%
  \end{alertblock}%
}
\newenvironment{paradoja}{%
  \begin{alertblock}{Exploratory Paradox}%
}{%
  \end{alertblock}%
}

\newenvironment{motivation}{%
  \begin{exampleblock}{Motivation: Data Science in Practice}%
}{%
  \end{exampleblock}%
}
\newenvironment{motivacion}{%
  \begin{exampleblock}{Motivation: Data Science in Practice}%
}{%
  \end{exampleblock}%
}

\newenvironment{quickchallenge}{%
  \begin{block}{Quick Team Challenge}%
}{%
  \end{block}%
}
\newenvironment{retoclase}{%
  \begin{block}{Quick Team Challenge}%
}{%
  \end{block}%
}

\title[Introduction to linear regression]{Section 09.02: Introduction to Linear Regression}\subtitle{Models, assumptions, and applications}\author[J. Castillo Colmenares]{Juliho Castillo Colmenares}\institute{Tecnológico de Monterrey}\date{}
\begin{document}
\begin{frame}[plain]\titlepage\end{frame}
\begin{frame}{Roadmap}\small\begin{enumerate}\item Context and models.\item Correlation, assumptions, and types.\item Applications and workflow.\end{enumerate}\end{frame}
\begin{frame}{Source and scope}\small The notes introduce linear regression as a predictive technique and prepare the chapter's mathematical and computational sections.\begin{block}{Guiding question}What makes a linear relationship a useful statistical model?\end{block}\end{frame}
\begin{frame}{Linear regression}\small Regression builds a directional relationship between explanatory variables and a response, using historical data to estimate parameters and produce predictions. It is a stochastic model: random error distinguishes observed data from a deterministic relationship.\end{frame}
\begin{frame}{Main uses}\small\begin{itemize}\item Forecast future values.\item Infer associations.\item Control processes.\item Conduct scientific research.\end{itemize}\end{frame}
\begin{frame}{Structural example}\small For house price $P$ with size $S$, amenities $A$, and transit access $T$,\[P=a_1S+a_2A+a_3T+\epsilon.\]\end{frame}
\begin{frame}{Parameters and variables}\small $P$ is the response; $S,A,T$ are known inputs; $a_1,a_2,a_3$ are parameters to estimate; $\epsilon$ collects unobserved factors.\end{frame}
\begin{frame}{Estimation}\small Estimate parameters from historical data. Once fitted, the model produces fitted values and new predictions.\end{frame}
\begin{frame}{Correlation and regression}\small\begin{itemize}\item Correlation: symmetric, unit-free linear association.\item Regression: directional relationship for modeling and prediction.\end{itemize}\end{frame}
\begin{frame}{Causal caution}\small A high correlation can suggest a linear fit, but it does not establish causality or replace an experimental design.\end{frame}
\begin{frame}{Linearity assumption}\small The conditional mean of the response should change linearly with predictors over the range of interest.\end{frame}
\begin{frame}{Independence assumption}\small Observations should be independent. Temporal, spatial, or clustered dependence requires a different model or design.\end{frame}
\begin{frame}{Homoscedasticity assumption}\small Error variance should be approximately constant across predictor levels.\end{frame}
\begin{frame}{Normality assumption}\small Normal errors matter especially for small-sample inference; normality is not required for every prediction task.\end{frame}
\begin{frame}{No multicollinearity}\small In multiple regression, highly correlated predictors make separate effects difficult to identify and coefficients unstable.\end{frame}
\begin{frame}{Simple regression}\small With one predictor,\[Y=\beta_0+\beta_1X+\epsilon.\]The slope summarizes expected change in $Y$ per unit of $X$.\end{frame}
\begin{frame}{Multiple regression}\small With several predictors,\[Y=\beta_0+\beta_1X_1+\cdots+\beta_pX_p+\epsilon.\]Each coefficient is a partial effect conditional on the others.\end{frame}
\begin{frame}{Least squares}\small Both forms estimate parameters by minimizing squared residuals. Variable selection must respect the question and assumptions.\end{frame}
\begin{frame}{Applications}\small Finance, marketing, medicine, engineering, and social science use regression to quantify associations and predict, always with validation.\end{frame}
\begin{frame}{Workflow}\small\begin{enumerate}\item Define response and predictors.\item Fit.\item Diagnose assumptions.\item Validate out of sample.\item Communicate uncertainty.\end{enumerate}\end{frame}
\begin{frame}{Synthesis}\small Linear regression is a stochastic model: it estimates a relationship, separates signal from noise, and requires explicit assumptions for interpretation.\end{frame}
\begin{frame}{Chapter connection}\small The next sections develop regression mathematics, implementations, validation, and extensions.\end{frame}
\end{document}
